\section{2. Related Work}

\subsection{Traditional Light Field Depth Estimation Methods}

Many depth estimation methods for light-field cameras have been proposed, primarily leveraging the geometric structures within Epipolar Plane Images (EPIs) \cite{EPIdepthestimationlightfie}. Wanner et al. \cite{WannerEPIstructuretensorgl} formulated depth estimation as an orientation analysis problem by computing the structure tensor $\mathbf{J}$ of EPIs. Specifically, they derived the disparity from the dominant eigenvector of $\mathbf{J}$, subsequently refining the initial depth map through a global variational optimization framework. This approach successfully demonstrates high accuracy and robustness on Lambertian surfaces with rich textures. However, the structure tensor relies strictly on the photo-consistency assumption, meaning that the linear structures in EPIs are explicitly assumed to be continuous and uniform across angular views. When encountering non-Lambertian materials, the specular reflections shift non-linearly, destroying the EPI linearity and causing the structure tensor to yield erroneous orientation estimates.

While EPI-based methods focus on angular geometric structures, defocus-based approaches exploit the focal stack to estimate depth by evaluating the sharpness of pixels across different focal planes \cite{defocusdepthestimationfocal}. To overcome the limitations of using either cue independently, Tao et al. \cite{Taodefocusparallaxfusionli} proposed a unified framework that seamlessly blends defocus and parallax (correspondence) cues. They defined a confidence measure for each cue and aggregated them to produce a remarkably robust depth map, effectively addressing textureless regions where parallax fails and occluded regions where defocus struggles. Moreover, their method significantly improves depth estimation in complex scenes by leveraging the complementary nature of these two cues. Unfortunately, both defocus and parallax cues inherently assume that the radiance of a 3D point $L(x,y,u,v)$ remains constant across different viewpoints and focal settings. Consequently, their fusion model fails to hold for specular or scattering surfaces, where the appearance changes dramatically with viewpoint and focus, leading to severe depth artifacts in non-Lambertian regions.

To address the violations of the Lambertian assumption, several heuristic strategies have been developed to explicitly handle specular highlights or non-Lambertian pixels \cite{specularhighlightremovalnon}. For instance, Tao et al. \cite{Taospecularremovallightfie} proposed a clustering-based method that identifies and eliminates specular pixels by analyzing the angular color distribution, subsequently inpainting the masked regions using neighboring Lambertian pixels. Similarly, other approaches \cite{heuristicmasklightfieldrob} employ hard thresholding on angular variance or epipolar consistency to mask out unreliable pixels before depth optimization. While these methods successfully mitigate the impact of strong specularities in controlled environments, they typically treat non-Lambertian reflections as outliers to be discarded or inpainted. In contrast to physical modeling, these heuristic rules lack the ability to decouple the underlying light transport mechanisms. Therefore, they fail if the scene contains mixed materials at the pixel level (e.g., urban scenes with complex scattering and specular coexistence), as hard masking inevitably discards valuable geometric information and introduces inpainting artifacts.

In short, traditional light field depth estimation methods encounter fundamental bottlenecks when applied to real-world, complex environments. First, they severely rely on the global Lambertian assumption; in non-Lambertian regions, the violation of view consistency leads to a complete failure of depth estimation, often exhibiting texture-copying artifacts. Second, conventional highlight processing primarily utilizes hard thresholds or heuristic rules, which cannot effectively handle complex pixels where multiple materials coexist in urban mixed scenes. Finally, these approaches lack explicit physical modeling of light-matter interactions, such as scattering and specular reflection, resulting in limited generalization capability under complex illumination conditions. In contrast, our model formulates a unified physical framework that operates on general conditions. By leveraging an MRI-like angular frequency analysis mechanism and a unified dual-mask physical model with inverse wavevector parsing, we decouple the mixed reflections and output continuous material weights rather than hard classifications. Furthermore, our component-aware adaptive depth fusion strategy thoroughly resolves the texture-copying failure, recovering accurate depth in addition to material properties, which has not been achieved by previous light field approaches.

\subsection{Deep Learning-based General Light Field Depth Estimation Methods}

To overcome the limitations of hand-crafted geometric cues, deep learning-based methods have been extensively explored to implicitly learn the mapping from 4D light field (LF) data to depth maps \cite{deeplearninglightfield4Dd}. These approaches primarily leverage convolutional neural networks (CNNs) to extract spatial and angular features from sub-aperture images (SAIs) or epipolar plane images (EPIs). As a pioneering work, EPINet \cite{EPINetEPICNNdepth} formulates depth estimation by feeding horizontal and vertical EPI slices into a multi-stream CNN architecture. By explicitly exploiting the linear structures within EPIs, EPINet successfully demonstrates remarkable accuracy on Lambertian surfaces. However, since it strictly relies on the photo-consistency assumption embedded in EPI linearity, its performance degrades dramatically when encountering non-Lambertian materials, where specular highlights shift non-linearly across views and destroy the expected EPI patterns.

While EPINet focuses on 1D EPI slices, subsequent methods extend the receptive field to 2D SAIs or 3D macropixels to capture richer spatial-angular correlations. LFNet \cite{LFNetmultiscalelightfield} proposes a multi-scale feature fusion strategy that aggregates local spatial details and global angular disparities by processing stacked SAIs through a hierarchical CNN. To further refine feature representation, DistgNet \cite{DistgNetdisentangledlightfi} explicitly decouples spatial and angular features using specialized disentangling modules and deformable convolutions. This decoupling mechanism allows the network to adaptively aggregate angular information, significantly improving depth estimation in occluded regions and texture-less areas. Despite these architectural advancements, both LFNet and DistgNet treat LF depth estimation as a pure pattern recognition task, implicitly assuming that the learned feature mappings are universally applicable to all surface reflectances without distinguishing the underlying physical reflection models.

Recently, to address the limited receptive field of CNNs, Transformer-based architectures have been introduced to model long-range spatial and angular dependencies in LF data \cite{LFTransformerattentionlight}. These methods typically employ self-attention mechanisms to compute global correlations across all SAIs, enabling the network to capture comprehensive scene geometry and effectively aggregate contextual information. Although Transformer-based models exhibit superior performance on standard benchmarks, their massive parameter spaces require extensive training data to converge. Unfortunately, the inherent scarcity of large-scale, accurately annotated non-Lambertian LF datasets severely restricts their generalization capability, making them prone to memorizing training artifacts rather than learning robust geometric representations in complex real-world scenarios.

In short, while existing deep learning-based general LF depth estimation methods have significantly advanced the state-of-the-art on standard benchmarks, they suffer from several critical limitations when applied to non-Lambertian and mixed-material scenes. First, as purely data-driven black-box models, they lack explicit physical priors of light reflection. In the presence of scarce non-Lambertian training data, these models are highly susceptible to overfitting, leading to poor generalization to complex real-world scenes. Second, directly introducing pre-trained large models (e.g., ResNet or Vision Transformers) as backbones exacerbates the severe overfitting problem on small-scale LF datasets, which fails to effectively enhance the depth accuracy in non-Lambertian regions. Most importantly, these learning-based methods cannot physically interpret the coexistence of multiple materials within mixed pixels. Lacking physical constraints, their depth predictions in non-Lambertian regions often degenerate into mere texture copying, failing to address the violation of the view-consistency assumption and lacking physical interpretability. In contrast, our approach departs from pure data-driven paradigms by introducing an angular frequency analysis mechanism and a unified dual-mask physical model, which explicitly resolves the multi-material coexistence and prevents texture-copying failures from a physical optics perspective.

\subsection{Deep Learning-based Physics-Aware and Non-Lambertian Depth Estimation Methods}

To mitigate the degradation caused by the violation of the Lambertian assumption, recent approaches have attempted to incorporate physical priors and material awareness into deep learning frameworks. A straightforward strategy is to employ learning-based classifiers to identify non-Lambertian regions before or during depth estimation. For instance, Wang et al. \cite{materialawarelightfieldspe} proposed a material-aware light field depth estimation network that utilizes a pre-trained 3-class CNN (categorizing pixels into diffuse, specular, and scattering) to generate a material mask. This mask subsequently guides a multi-branch network to apply tailored disparity estimation strategies for different surface types. While this method remarkably improves robustness in purely specular regions, it strictly relies on the accuracy of the local patch-based classifier, which is highly sensitive to the limited spatial context.

While the aforementioned method treats material classification as a separate pre-processing step, other works explicitly embed physical rendering models into the network architecture to jointly estimate depth and material properties. Chen et al. \cite{physicsbasedlightfieldinve} formulated a physics-aware inverse rendering framework that integrates the Bidirectional Reflectance Distribution Function (BRDF) into a differentiable pipeline. By minimizing the photometric error between the rendered and observed sub-aperture images, their model successfully decouples depth and spatially-varying BRDF parameters. Although this physically grounded approach demonstrates strong interpretability on synthetic datasets, it typically assumes a single dominant material per pixel. Unfortunately, this hard classification (either-or) paradigm fails to output continuous material weights, making it fundamentally incapable of addressing the pixel-level material mixing (Mixed) prevalent in real-world scenes.

In contrast to joint inverse rendering, another line of research focuses on explicitly separating reflection components to restore the violated photo-consistency. Zhang et al. \cite{specularseparationlightfiel} developed a specular-aware EPI refinement network that first predicts the non-linear shift of specular highlights and then separates the diffuse and specular components. The depth is subsequently estimated primarily from the restored diffuse EPIs. This component-decoupling strategy significantly alleviates the "texture copying" artifact caused by shifting highlights. However, their physical model is predominantly designed for ideal mirror-like specular reflections and lacks a unified theoretical framework to simultaneously explain complex scattering or translucent materials. Moreover, the iterative optimization required for component separation incurs high computational complexity, hindering lightweight, real-time pixel-level parsing.

In short, despite the progress made by these physics-aware and non-Lambertian methods, they collectively exhibit three critical limitations. First, existing learning-based material classifiers (e.g., 3-class CNNs) suffer from severe information deficiency under the small receptive fields of light field patches (e.g., $9 \times 9$ patches), leading to extremely low classification accuracy and failing to provide reliable priors for depth estimation. Second, current physical models predominantly adopt hard classification mechanisms, which cannot output continuous material weights and thus struggle to handle the intricate pixel-level multi-material mixing in real scenarios. Third, there is a lack of a unified physical theoretical framework to simultaneously explain diffuse, specular, and scattering reflections, resulting in high computational complexity and poor generalization. In contrast, our proposed method introduces an MRI-like angular frequency analysis mechanism for lightweight, highly accurate material classification without massive data training. Furthermore, we establish a unified dual-mask physical model with reverse wavevector parsing that outputs continuous material weights, effectively resolving the mixed-pixel dilemma. By formulating this unified framework, our approach avoids the heavy computational overhead of iterative rendering and enables a component-aware adaptive depth fusion strategy, thoroughly resolving the "texture copying" failure caused by the violation of photo-consistency in non-Lambertian regions.

While the aforementioned learning-based and physically-guided methods have advanced non-Lambertian light field depth estimation, they generally struggle with a heavy reliance on massive training data, inflexible hard material classification in mixed pixels, and severe depth degradation (e.g., texture copying) when photo-consistency assumptions are violated. Inspired by these observations, we propose a unified dual-mask physical modeling approach that addresses the aforementioned limitations by integrating an MRI-inspired angular frequency mechanism for lightweight material classification and a reverse wavevector analysis to yield continuous material weights. Furthermore, a component-aware adaptive fusion strategy is incorporated to fundamentally resolve the texture-copying failure, enabling robust and physically interpretable depth estimation across pure diffuse, non-Lambertian, and complex mixed scenes.